\documentclass[12pt,a4paper]{article}
\usepackage[a4paper,margin=1.20in]{geometry}
\usepackage[T1]{fontenc}
\usepackage[utf8]{inputenc}
\usepackage{lmodern}
\usepackage{microtype}
\usepackage{amsmath,amssymb,amsthm,mathtools}
\usepackage{enumitem}
\usepackage[hidelinks]{hyperref}
\usepackage{titlesec}
\titleformat{\section}{\large\bfseries}{\thesection}{0.6em}{}
\titleformat{\subsection}{\normalsize\bfseries}{\thesubsection}{0.6em}{}
\newcommand{\webersection}[2]{\clearpage\section*{\S #1. #2}}
\newcommand{\booktitle}[1]{\clearpage\section*{#1}}
\newcommand{\OmegaK}{\Omega}
\newcommand{\Om}{\Omega_m}
\newcommand{\Hm}{H_m}
\newcommand{\zetaS}{\zeta}
\newcommand{\idealp}{\mathfrak p}
\newcommand{\idealq}{\mathfrak q}
\newcommand{\ideala}{\mathfrak a}
\newcommand{\idealb}{\mathfrak b}
\newcommand{\C}{\mathcal C}
\newcommand{\N}{\operatorname{N}}
\newcommand{\Tr}{\operatorname{Tr}}
\newcommand{\Gal}{\operatorname{Gal}}
\newcommand{\ord}{\operatorname{ord}}
\newcommand{\Li}{\operatorname{Li}}
\newcommand{\Res}{\operatorname{Res}}
\setlength{\parindent}{0pt}
\setlength{\parskip}{0.98em}
\linespread{1.84}
\title{Heinrich Weber, \emph{Lehrbuch der Algebra}, Volume II\\English Translation: Sections 176--194\\Source pp. 641--710}
\author{}
\date{}
\begin{document}
\maketitle

\booktitle{Twenty-first Section. Cyclotomic Fields}

\webersection{176}{The Real Field Contained in \(\Om\)}

Every number of the cyclotomic field \(\Om\) passes, under the substitution
\[
  \zetaS\longmapsto \zetaS^{-1},
\]
into its conjugate imaginary value; this is the same operation as replacing \(i\) by \(-i\) in the complex representation.  The product of two such conjugate imaginary numbers is the square of the absolute value of either of them.

A number of \(\Om\) is real if and only if it remains unchanged by the interchange \(\zetaS\mapsto \zetaS^{-1}\).  Every such number is a rational function of the two-term period
\[
  \rho=\zetaS+\zetaS^{-1},
\]
and it belongs to a real field of degree \(\frac12\varphi(m)\), a divisor of \(\Om\).  This field will be called the real field contained in \(\Om\), and will be denoted by \(\Hm\).  This terminology is not meant to say that it is the only real divisor of \(\Om\); it is the distinguished real subfield fixed by complex conjugation.

Let
\[
  \mu=\varphi(m),\qquad [\Hm:\mathbb Q]=\frac12\mu.
\]
Then the numbers
\[
  1,
  \quad \zetaS+\zetaS^{-1},
  \quad \zetaS^2+\zetaS^{-2},
  \quad \ldots,
  \quad \zetaS^{\mu/2-1}+\zetaS^{-\mu/2+1}
\tag{1}
\]
form a basis of the whole numbers of \(\Hm\).  Indeed, by the earlier basis theorem for \(\Om\), a whole number of \(\Om\) has the form
\[
  \omega=x_0+x_1\zetaS+x_2\zetaS^2+\cdots+x_{\mu-1}\zetaS^{\mu-1},
\tag{2}
\]
where the coefficients \(x_i\) are rational integers.  The condition that \(\omega\) lie in \(\Hm\) is that it be unchanged by \(\zetaS\mapsto\zetaS^{-1}\).  Subtracting the transformed expression from (2), dividing by \(\zetaS-\zetaS^{-1}\), and using the irreducibility of the cyclotomic equation, one obtains the relations that pair the coefficients belonging to reciprocal exponents.  Thus every whole number of \(\Hm\) can be written in the form
\[
  \omega=x_0+x_1(\zetaS+\zetaS^{-1})+x_2(\zetaS^2+\zetaS^{-2})+\cdots.
\tag{3}
\]
This proves the assertion.

Equivalently, one may choose the powers of
\[
  \rho=\zetaS+\zetaS^{-1}
\]
as basis.  The system
\[
  1,\rho,\rho^2,\ldots,\rho^{\mu/2-1}
\tag{4}
\]
is obtained integrally from (1), and conversely the expressions (1) are integral polynomials in \(\rho\).  The number \(\rho\) is the root of an irreducible equation of degree \(\mu/2\), which we write
\[
  F_m(\rho)=0.
\]
If \(f_m(x)=0\) denotes the cyclotomic equation for \(\zetaS\), then the identity
\[
  f_m(x)=x^{\mu/2}F_m(x+x^{-1})
\tag{5}
\]
connects the two equations.  Differentiating gives
\[
  f_m'(\zetaS)=\zetaS^{\mu/2}F_m'(\rho)(1-\zetaS^{-2}).
\tag{6}
\]

This formula gives the fundamental number of \(\Hm\).  Since \(\Hm\) is real, its discriminant is positive.  If \(\N_1\) denotes the norm formed in the real field \(\Hm\), while \(\N\) denotes the norm in the full cyclotomic field \(\Om\), then the norm in \(\Om\) of a number belonging to \(\Hm\) is the square of its norm in \(\Hm\).  Hence (6), together with the formulae already found for \(f_m'(\zetaS)\) and for \(\N(1-\zetaS^{-2})\), yields that the discriminant \(\Delta_1\) of \(\Hm\) is always a power of the prime \(q\) when \(m=q^a\).

More explicitly, in Weber's notation one obtains
\[
  \Delta_1=q^{\frac12(q-1)q^{a-1}-q^{a-1}}
  \quad(q\text{ odd}),
\tag{7a}
\]
and for \(q=2\) a corresponding power of \(2\), obtained from the special value
\[
  \N(1-\zetaS^{-2})=4.
\tag{7b}
\]
The important point for what follows is not the precise exponent, but the fact that no prime except \(q\) divides the discriminant of the real subfield.  The ramification in \(\Hm\) is therefore concentrated above the same rational prime as in \(\Om\), though with half the degree forced by passage to the fixed field of complex conjugation.

\webersection{177}{The Prime Ideals of the Field \(\Hm\)}

We now have to determine how the prime ideals of the real field \(\Hm\) stand in relation to those of the full cyclotomic field \(\Om\).

Let \(\idealp\) be a prime ideal of \(\Hm\), of degree \(h\), lying above a rational prime \(p\).  A prime ideal of \(\Om\) above \(\idealp\) can be ramified only if \(p=q\), because in \(\Om\) all primes except \(q\) are unramified.  If \(\idealp\) lies above \(q\), then it must be a power of the prime divisor
\[
  \lambda=1-\zetaS
\]
of \(q\) in \(\Om\).  Writing
\[
  \idealp=\lambda^\ell,
\]
and forming norms, the norm from \(\Om\) is the square of the norm from \(\Hm\).  Thus
\[
  q^{2h}=q^\ell.
\]
It follows that \(\ell=2h\).  But \((1-\zetaS)(1-\zetaS^{-1})\) is a whole number of \(\Hm\), and is associated with \(\lambda^2\).  Therefore \(\ell\) cannot exceed \(2\).  Hence \(h=1\) and \(\ell=2\).  We get:

\textbf{1.} The rational prime \(q\) is the \(\frac12\mu\)-th power of a prime number of first degree in \(\Hm\).

Now let \(p\ne q\), and let \(\mathfrak P\) be a prime factor of \(p\) in \(\Om\), of degree \(f\).  Under the substitution \(\zetaS\mapsto\zetaS^{-1}\), the ideal \(\mathfrak P\) goes to another prime ideal \(\mathfrak P'\).  The prime ideal \(\idealp\) of \(\Hm\) below it is divisible by both \(\mathfrak P\) and \(\mathfrak P'\).  On the other hand, the product \(\mathfrak P\mathfrak P'\) is invariant under conjugation and hence belongs to \(\Hm\).

There are two cases.  If \(\mathfrak P\ne\mathfrak P'\), then
\[
  \idealp=\mathfrak P\mathfrak P',
\]
and the degree of \(\idealp\) is \(h=f\).  If, however, \(\mathfrak P=\mathfrak P'\), then \(\idealp=\mathfrak P\), and norm formation gives
\[
  2h=f.
\]
Which case occurs depends on whether the prime ideal \(\mathfrak P\) is invariant under complex conjugation.  In the group of substitutions of \(\Om\), this means that the substitution \(\zetaS\mapsto\zetaS^{-1}\) belongs to the decomposition group of \(\mathfrak P\).  Equivalently, there must be an exponent \(h\) for which
\[
  p^h\equiv -1\pmod m.
\tag{1}
\]
If no such exponent exists, then the two conjugate prime ideals of \(\Om\) remain distinct over \(\Hm\).

The result may be stated as follows.

\textbf{2.} The rational primes \(p\ne q\) divide into two kinds.

The first kind consists of those primes for which
\[
  p^h\equiv -1\pmod m
\]
for some positive exponent \(h\).  If such a prime belongs to the exponent \(f\) modulo \(m\), and if
\[
  \mu=ef,
\]
then it decomposes in \(\Hm\) into \(e\) prime factors of degree \(f/2\).  These prime factors are not decomposed further in \(\Om\).

The second kind consists of primes for which no power is congruent to \(-1\) modulo \(m\).  Such a prime decomposes in \(\Hm\) into \(\frac12e\) prime factors of degree \(f\); each of these prime factors splits in \(\Om\) into two conjugate factors.

For primes of the first kind the exponent \(f\) must be even.  If \(m\) is odd, the converse also holds: every prime belonging to an even exponent is of the first kind.  For if \(p^f\equiv1\pmod m\) and \(f\) is even, then
\[
  (p^{f/2}-1)(p^{f/2}+1)\equiv0\pmod m.
\]
The first factor cannot be divisible by \(m\), since otherwise the exponent would be smaller than \(f\).  As \(m
e2\), the two factors cannot both have the same prime divisor in such a way as to spoil the conclusion; therefore
\[
  p^{f/2}\equiv -1\pmod m.
\]

If \(m\) is a power of \(2\), the congruence \(p^h\equiv-1\pmod m\) is possible only when already
\[
  p\equiv -1\pmod m.
\]
Indeed every even power of an odd number is congruent to \(+1\) modulo \(8\), and if \(h\) is odd then \(p^h+1\) is divisible by no higher power of \(2\) than \(p+1\).

Thus:

\textbf{3.} If \(m\) is odd, then the primes of the first kind are precisely those with even exponent, and those of the second kind have odd exponent.  If \(m\) is even, then primes of the form \(hm-1\) are of the first kind; all other odd primes are of the second kind.

\webersection{178}{The Units of the Field \(\Hm\)}

The units belonging to \(\Hm\) are especially important in the applications.  The units of the whole cyclotomic field \(\Om\) may be reduced to them by the following theorem.

\textbf{1.} Every unit of \(\Om\) is the product of a root of unity and a real unit of \(\Hm\).

Let \(E(\zetaS)\) be any unit of \(\Om\).  The number \(E(\zetaS^{-1})\) is its conjugate imaginary value and is again a unit.  Therefore the quotient
\[
  \frac{E(\zetaS)}{E(\zetaS^{-1})}
\]
is a unit all of whose conjugates have absolute value \(1\).  By the theorem on units of absolute value \(1\), it is a root of unity.  Hence
\[
  E(\zetaS)=\pm \zetaS^h E(\zetaS^{-1})
\tag{1}
\]
with an integral exponent \(h\).

If \(m\) is odd, we may suppose that \(h\) is even, after replacing \(h\) by \(h+m\) if necessary.  Moreover only the plus sign in (1) is admissible.  For if the minus sign occurred, then after putting \(\zetaS=1\) one would get
\[
  E(1)=-E(1),
\]
so that \(2E(1)\equiv0\pmod{1-\zetaS}\).  Since \(1-\zetaS\) divides \(q\) and is prime to \(2\), this would force \(E(\zetaS)\) to be divisible by \(1-\zetaS\), contrary to the assumption that it is a unit.  Consequently
\[
  \zetaS^{-h/2}E(\zetaS)
\]
is fixed by complex conjugation.  It is therefore a real unit; if it is denoted by \(\varepsilon(\zetaS)\), then
\[
  E(\zetaS)=\zetaS^{h/2}\varepsilon(\zetaS).
\tag{2}
\]

If \(m\) is a power of \(2\), then \(\zetaS^{m/2}=-1\), and with \(\zetaS\) also \(-\zetaS\) is a primitive \(m\)-th root of unity.  Replacing \(h\) by \(h+m/2\) if necessary, we may again arrange that the plus sign occurs in (1).  It remains only to show that \(h\) is even.  Write
\[
  E(\zetaS)=\sum_s x_s\zetaS^s,
\tag{3}
\]
with rational integral coefficients \(x_s\).  Formula (1) relates the coefficient \(x_s\) to a coefficient \(x_{s'}\) with
\[
  s+s'\equiv h\pmod m.
\]
If \(h\) were odd, then one of the two exponents would be even and the other odd; the corresponding terms would group into expressions divisible by \(1-\zetaS\).  Then the whole unit \(E(\zetaS)\) would be divisible by \(1-\zetaS\), impossible.  Hence \(h\) is even, and (2) holds here also.  The theorem is proved.

We now point out a useful family of real units.  Since all the numbers \(1-\zetaS^n\), with \((n,m)=1\), are associated, the quotient
\[
  \frac{1-\zetaS^n}{1-\zetaS^{n'}}
\tag{4}
\]
is a unit whenever \(n,n'\) are prime to \(m\).  Multiplying numerator and denominator by suitable powers of \(\zetaS\), this quotient may be written in the real form
\[
  \frac{\zetaS^{n/2}-\zetaS^{-n/2}}{\zetaS^{n'/2}-\zetaS^{-n'/2}}
  =\frac{\sin(n\pi/m)}{\sin(n'\pi/m)}.
\tag{5}
\]
Thus (5) gives real units in \(\Hm\).  If \(m\) is even, one may put \(n'=n+m/2\), and obtains the particularly simple units
\[
  \tau_n=\tan\frac{n\pi}{m},
\tag{6}
\]
where \(n\) is odd.  Replacing \(n\) by \(m/2-n\) changes \(\tau_n\) into its reciprocal.  If \(n\) runs through
\[
  1,3,5,\ldots,\frac m2-1,
\]
one obtains positive proper fractional values, a system that later reappears in the formulae for the class number of real cyclotomic fields.

\booktitle{Twenty-second Section. Abelian Fields and Cyclotomic Fields}

\webersection{179}{Simple Abelian Fields}

We now turn to one of the most interesting applications of the theory of algebraic numbers.  It is the proof of the general theorem:

\textbf{I.} Every number field which is Abelian over the absolute rational domain is a cyclotomic field.

Once this theorem has been proved, the investigations of the earlier section on cyclotomic fields acquire a wider meaning: the method given there does not merely construct all cyclotomic fields, but all Abelian number fields.

Let
\[
  \Omega=\mathbb Q(x)
\]
be an Abelian field of degree \(m\).  Thus \(x\) is a root of an irreducible Abelian equation of degree \(m\), the rational numbers being the ground domain.  The Galois group \(G\) of the equation is Abelian and has the same degree \(m\) as the field; it is also the group of the field.

If \(x_1\) is a non-primitive number of \(\Omega\), then \(\Omega_1=\mathbb Q(x_1)\) is again an Abelian field.  We call it a proper divisor of \(\Omega\), and its degree is a proper divisor of the degree of \(\Omega\).  For if \(A\) is the group to which \(x_1\) belongs, then the quotient of \(G\) by \(A\) is the group of \(\Omega_1\), and this group is Abelian.

If there are two non-primitive numbers \(x_1,x_2\) in \(\Omega\) for which
\[
  \Omega=\mathbb Q(x_1,x_2),
\tag{1}
\]
then the fields \(\Omega_1=\mathbb Q(x_1)\), \(\Omega_2=\mathbb Q(x_2)\) are proper divisors of \(\Omega\), and \(\Omega\) is said to be composed from them.  If no such representation is possible, the field is called simple.

Thus, if the theorem has already been proved for \(\Omega_1\) and \(\Omega_2\), it also follows for a field composed from them: every number of \(\Omega_1\) and \(\Omega_2\) is rationally expressible by roots of unity, hence every number of their compositum is so expressible.  By repeating the decomposition, and since the degrees strictly decrease, it is enough to prove theorem I for simple Abelian fields.

A criterion for simplicity is read from the group.  Suppose the Abelian group \(G\) is decomposable into two proper components \(A,B\), meaning that every element of \(G\) is represented in one and only one way as a product of an element of \(A\) and an element of \(B\):
\[
  G=AB.
\tag{2}
\]
Then \(|G|=|A|\,|B|\).

\textbf{1.} If the group of an Abelian field is decomposable, then the field is composed.

Indeed, let \(|G|=m\), \(|A|=a\), \(|B|=b\).  Choose a number \(\xi\) belonging to the group \(B\), so that it has \(a\) conjugate values, and a number \(\eta\) belonging to the group \(A\), so that it has \(b\) conjugate values.  With suitable rational numbers \(\alpha,\beta\), the number
\[
  x=\alpha\xi+\beta\eta
\]
has all \(m=ab\) conjugates distinct.  Hence \(\Omega=\mathbb Q(x)\) is generated by \(\mathbb Q(\xi)\) and \(\mathbb Q(\eta)\), as claimed.

Using the representation of Abelian groups by a basis, one concludes:

\textbf{2.} Every Abelian field can be composed from fields whose degrees are prime powers and whose groups are cyclic.  The degrees of these composing fields are the invariants of the group of the original field.

This is supplemented by the converse statement needed for the reduction:

\textbf{3.} An Abelian field of prime-power degree with cyclic group is simple.

Let \(G\) be cyclic of order \(m=q^a\).  Its subgroups are nested; any two proper subgroups are contained in a larger proper subgroup.  If the field could be composed from two proper divisor fields, the corresponding subgroups would produce a decomposition of the cyclic group into two proper components.  This is impossible in a cyclic group of prime-power order.  More explicitly, if the subgroups corresponding to two divisor fields are proper, then the field obtained from both is still a proper divisor.  Thus the whole field cannot be produced from them.

It remains, therefore, to prove theorem I for a cyclic Abelian extension of degree \(q^a\), where \(q\) is prime.  This is the purpose of the following paragraphs.

\webersection{180}{The Resolvents}

We may now confine ourselves to Abelian fields whose degree is a prime power and whose group is cyclic.  For the intended application it was therefore enough, in the preceding section, to study those cyclotomic fields \(\Om\) for which
\[
  m=q^a
\tag{1}
\]
is a prime power.  In the case \(q=2\) we assume \(a>1\).  Let \(n\) denote a number not divisible by \(q\).  Let \(\zetaS\) be a primitive \(m\)-th root of unity, and let \(\Om\) be the field of rational functions of \(\zetaS\).

Let
\[
  \Omega=\mathbb Q(x_0)
\]
be a simple Abelian field of degree \(m\), with cyclic Galois group.  Write its roots in cyclic order
\[
  x_0,x_1,\ldots,x_{m-1},
\]
where the generating substitution sends
\[
  x_0\mapsto x_1\mapsto x_2\mapsto\cdots\mapsto x_{m-1}\mapsto x_0.
\tag{2}
\]
Every cyclic function of these roots is rational.  We have to show that the roots \(x_i\) are rational functions of roots of unity.

Adjoin the root of unity \(\zetaS\).  The composed field
\[
  \mathbb Q(x_0,\zetaS)
\]
contains both \(\Omega\) and \(\Om\).  A number of this field which remains fixed by all cyclic substitutions of the \(x_i\) has coefficients, in its expression as a function of \(\zetaS\), which are symmetric functions of the roots \(x_i\), and hence rational.  It therefore lies in \(\Om\).  This remains true even if the equation for \(x_0\) becomes reducible after adjoining \(\zetaS\).

We apply this to the Lagrange resolvents
\[
  (\zetaS,x_0)=x_0+\zetaS x_1+\zetaS^2x_2+\cdots+\zetaS^{m-1}x_{m-1},
\tag{3}
\]
and, more generally,
\[
  (\zetaS^s,x_0)=x_0+\zetaS^s x_1+\zetaS^{2s}x_2+\cdots+\zetaS^{(m-1)s}x_{m-1}.
\tag{4}
\]
These resolvents determine \(x_0\), for by summing over all residues \(s\) modulo \(m\),
\[
  m x_0=\sum_{s=0}^{m-1}(\zetaS^s,x_0).
\tag{5}
\]
Thus it is enough to prove that each resolvent \((\zetaS^s,x_0)\) is a cyclotomic number.

If
\[
  (\zetaS^s,x_j)=x_j+\zetaS^s x_{j+1}+\cdots+\zetaS^{(m-1)s}x_{j+m-1}
\]
(indices modulo \(m\)), then the cyclic substitution gives the fundamental relation
\[
  \zetaS^{sj}(\zetaS^s,x_j)=(\zetaS^s,x_0).
\tag{6}
\]
It follows in particular that the \(m\)-th power
\[
  \omega=(\zetaS,x_0)^m
\tag{7}
\]
is fixed by the cyclic substitutions of the \(x_i\), and hence belongs to \(\Om\).  Thus the prime factorization of \(\omega\) in the cyclotomic field can be studied by the ideal theory already developed.

If \(\omega_n\) denotes the conjugate of \(\omega\) obtained by \(\zetaS\mapsto\zetaS^n\), then the relation among the resolvents shows that the system \(\omega_n\) satisfies strong congruence and factorization conditions.  The next paragraph derives the form of these conditions; the two following paragraphs prove the two auxiliary assertions needed to complete the proof of theorem I.

\webersection{181}{Preparation for the Proof}

Let
\[
  \omega=(\zetaS,x_0)^m
\]
be the \(m\)-th power of the principal resolvent.  Since \(\omega\) belongs to \(\Om\), it has in \(\Om\) a factorization into prime functionals.  We shall separate the part depending on the primes above \(q\) from the part depending on the other primes.

The resolvents satisfy multiplication formulae that, in the earlier theory of cyclotomy, appeared in the special form of products of periods.  Applied here, they show that the prime factors of \(\omega\) outside \(q\) occur in powers whose exponents are controlled by the conjugates of the resolvent.  If a rational prime \(p\ne q\) divides the norm of the resolvent, then the exponents of the prime ideals above \(p\) may be adjusted by multiplying by suitable cyclotomic resolvents of the form
\[
  (\zetaS^{mn},\eta)^m,
\]
where \(\eta\) is a root belonging to a higher cyclotomic field.  In this way one obtains a factorization
\[
  \omega_n=\varepsilon_n\varphi_n^m
       \prod_p (\zetaS^{m n_p},\eta_p)^{m e_p},
\tag{1}
\]
where \(\varepsilon_n\) is a unit, \(\varphi_n\) is a functional of \(\Om\), and the product extends over rational primes \(p\equiv1\pmod q\), with repetitions if necessary.

The essential point of (1) is that all factors except \(\varepsilon_n\) and \(\varphi_n\) are known cyclotomic numbers.  Moreover, the conjugate functionals \(\varphi_n\) and units \(\varepsilon_n\) satisfy two characteristic properties.

\textbf{7.} The functionals \(\varphi_n\) are associated with numbers of the field \(\Om\); in other words, they belong to the principal class.

This is first known only for certain products of the \(\varphi_n\).  From (1) and from the fact that \(\omega_n\omega_{-n}\) is an \(m\)-th power in \(\Om\), it follows that products of the form
\[
  \varphi_n\varphi_{-n}
\]
are principal; more generally the products obtained by conjugation are principal.  The first auxiliary theorem will show that this already implies that each \(\varphi_n\) itself is principal.

\textbf{8.} The conjugate functionals \(\varphi_n\) have the property that all products
\[
  \varphi_n\varphi_{-n},\quad \varphi_n\varphi_{n'},\quad\ldots
\]
which arise from the resolvent relations lie in the principal class.

\textbf{9.} The units \(\varepsilon_n\) have the property that the corresponding products
\[
  \varepsilon_n\varepsilon_{-n},\quad \varepsilon_n\varepsilon_{n'},\quad\ldots
\]
are \(m\)-th powers of units.

Thus the proof of theorem I is reduced to two lemmas.

\textbf{A.} If the conjugate functionals \(\varphi_n\) satisfy the product conditions just described, then each \(\varphi_n\) belongs to the principal class.

Assuming A, choose numbers \(\alpha_n\in\Om\) associated with \(\varphi_n\).  The factorization (1) becomes, after absorbing principal factors,
\[
  \omega_n=\varepsilon_n\alpha_n^m
       \prod_p (\zetaS^{m n_p},\eta_p)^{m e_p}.
\tag{2}
\]
Only the units \(\varepsilon_n\) remain to be dealt with.

\textbf{B.} If a system of conjugate numerical units \(\varepsilon_n\) satisfies the product condition of theorem 9, then each \(\varepsilon_n\) is, apart from a root of unity \(\zetaS^a\), an \(m\)-th power of a unit of \(\Om\).

Assuming B, (2) shows that each \(\omega_n\) is an \(m\)-th power multiplied by known cyclotomic factors.  Hence the resolvent \((\zetaS^n,x_0)\) itself lies in a cyclotomic field.  Formula (5) of the preceding paragraph then gives \(x_0\) as a rational function of roots of unity.  This proves that \(\Omega\) is a cyclotomic field.

The next two paragraphs prove A and B when \(m\) is odd.  The even case requires a supplementary reduction and is deferred to the following discussion.

\webersection{182}{Proof of the First Auxiliary Lemma for Odd \(m\)}

We assume here that \(m=q^a\) is odd.  The first auxiliary lemma may be formulated as follows.

Let \(\varphi_n\) be a system of conjugate functionals in the cyclotomic field \(\Om\), where \(n\) runs through the residues prime to \(m\).  Suppose that all products occurring in the resolvent relations are principal.  Then each \(\varphi_n\) itself is principal.

It is enough to use the class group.  Let \([\varphi]\) be the class of \(\varphi_1\), and let \([\varphi_n]\) be its conjugates under \(\zetaS\mapsto\zetaS^n\).  The hypotheses say that certain products of these conjugate classes are the principal class.  Passing to exponents in the finite Abelian group of ideal classes, one obtains congruences for the exponents of the conjugate classes.

The decisive arithmetic ingredient is a congruence for a sum of integral parts.  If \(t\) runs through a complete reduced residue system modulo \(m\), let \(t'\) be defined by
\[
  tt'\equiv1\pmod m.
\]
For \(n\) prime to \(m\) define
\[
  S_n=\sum_t t'\left\lfloor\frac{nt}{m}\right\rfloor.
\tag{1}
\]
The class relations imply that \([\varphi]^ {S_n}\) is principal.  Hence if for some \(n\) the integer \(S_n\) is not divisible by \(q\), the class \([\varphi]\) is killed by an exponent prime to \(q\).  Since the class group contribution here is a \(q\)-group, this forces \([\varphi]
=1\).

We therefore have only to show that, among the numbers \(S_n\), at least one is not divisible by \(q\).  Weber reduces this to the case in which \(m\) itself is prime.  Write
\[
  m=qm',
\]
and decompose
\[
  t=t_1+qt_2,
\]
where \(t_1=1,
2,\ldots,q-1\) and \(t_2=0,1,\ldots,m'-1\).  Choose \(t_1'\) from
\[
  t_1t_1'\equiv1\pmod q.
\]
Then, modulo \(q\), the reciprocal \(t'\) is congruent to \(t_1'\), and the summation in (1) separates into a part depending on \(m'\) and a part depending only on \(q\).  The elementary identity for integral parts is
\[
  \left\lfloor\frac{n(t_1+qt_2)}{qm'}\right\rfloor
  =\left\lfloor\frac{nt_2}{m'}\right\rfloor
\tag{2}
\]
or that value plus \(1\), according as a simple inequality involving the fractional part is or is not satisfied.  Counting the occurrences of the second case gives
\[
  S_n\equiv n\sum_{t_1=1}^{q-1} t_1'\left\lfloor\frac{nt_1}{q}\right\rfloor\pmod q.
\tag{3}
\]
Thus the question is reduced to the same question for the prime modulus \(q\).

For a prime modulus, compare \(S_n\) and \(S_{q-n}\).  The two corresponding fractional parts add to \(1\), and after multiplication by \(t_1'
\) and summation one obtains a congruence of the form
\[
  S_n+S_{q-n}\not\equiv0\pmod q.
\tag{4}
\]
Therefore at least one of \(S_n,S_{q-n}\) is not divisible by \(q\).  This proves the existence of the required exponent and hence proves the first auxiliary lemma for odd \(m\).

The proof is characteristic of Weber's treatment: the ideal-class assertion is converted into a finite congruence in the group of residues, and the group-theoretic obstruction is removed by an explicit calculation with integral parts.

\webersection{183}{Proof of the Second Auxiliary Lemma for Odd \(m\)}

The second lemma, still for odd \(m\), is the following.

\textbf{2.} Let \(\varepsilon_n\) be a system of conjugate numerical units of \(\Om\) satisfying the condition that
\[
  \varepsilon_n\varepsilon_{-n}
\tag{1}
\]
is an \(m\)-th power of a unit of \(\Om\) for every \(n\) prime to \(m\).  Then each \(\varepsilon_n\) itself is an \(m\)-th power of a unit, multiplied by a root of unity of the form \(\zetaS^a\).

By the theorem of \S178, every unit of \(\Om\) can be written as a root of unity times a real unit of \(\Hm\).  Thus
\[
  \varepsilon=\pm\zetaS^a E(\zetaS),
\qquad
  \varepsilon_n=\pm\zetaS^{an}E(\zetaS^n),
\tag{2}
\]
where \(E(\zetaS)\) is real, i.e.
\[
  E(\zetaS)=E(\zetaS^{-1}).
\tag{3}
\]
For \(n=-1\), condition (1) gives
\[
  \varepsilon_1\varepsilon_{-1}=E(\zetaS)^2
\]
up to a root of unity, and in fact the real part is an \(m\)-th power of a unit.  Hence
\[
  E(\zetaS)^2=U(\zetaS)^m
\tag{4}
\]
for some unit \(U(\zetaS)\) of \(\Om\).

Since \(m\) is odd, one can choose integers \(\lambda,\mu\) such that
\[
  2\lambda+m\mu=1.
\tag{5}
\]
Raising (4) to the power \(\lambda\) and multiplying by a suitable \(m\)-th power, we obtain
\[
  E(\zetaS)=V(\zetaS)^m.
\tag{6}
\]
Substitution \(\zetaS\mapsto\zetaS^n\) gives the same conclusion for every conjugate \(E(\zetaS^n)\).  Therefore each \(\varepsilon_n\) is an \(m\)-th power of a unit multiplied by a power of \(\zetaS\), which is exactly the assertion.

Combining \S182 and \S183 proves theorem I for cyclic Abelian fields of odd prime-power degree.  Together with the reduction of \S179, it proves the Kronecker-Weber theorem for Abelian fields whose cyclic prime-power constituents all have odd degree.

For even degree the argument just used fails precisely at (5).  The equation \(2\lambda+m\mu=1\) is no longer solvable.  The remaining difficulty is therefore a genuine two-adic one, and Weber treats it by a preliminary separation of the even case.

\webersection{184}{Preliminaries on the Case of Even \(m\)}

When \(m\) is even, the preceding proof of the second auxiliary lemma cannot be carried over unchanged.  The obstruction is elementary but decisive: from the fact that the square of a real unit is an \(m\)-th power one can no longer conclude that the real unit itself is an \(m\)-th power, because \(2\) is not invertible modulo \(m\).

The remedy is to isolate the even part.  Let \(m=2^a\) with \(a>1\), and let \(\zetaS\) be a primitive \(m\)-th root of unity.  Then \(-\zetaS\) is again a primitive \(m\)-th root if necessary after changing exponents, and complex conjugation is still \(\zetaS\mapsto\zetaS^{-1}\).  The unit theorem of \S178 remains available, but the root of unity attached to a unit must be chosen with more care.  In particular, signs and powers of \(\zetaS^{m/2}=-1\) cannot be suppressed by the same argument used in the odd case.

The proof begins by observing that the only additional ambiguity comes from signs.  If a conjugate system of units \(\varepsilon_n\) satisfies
\[
  \varepsilon_n\varepsilon_{-n}=U_n^m,
\tag{1}
\]
then after extracting the real factors as in \S178 one obtains real units \(E(\zetaS^n)\) for which the products
\[
  E(\zetaS^n)E(\zetaS^{-n})=E(\zetaS^n)^2
\]
are \(m\)-th powers up to a sign.  Since \(m\) is a power of \(2\), the sign must be absorbed by adjoining, if necessary, a fourth root of unity.  This reduces the assertion to a field of twice the conductor.  In modern language, the proof passes from a cyclic two-power extension to a cyclotomic field with conductor enlarged by the minimal two-power needed to make the unit equation soluble.

Weber's conclusion is that the even case also reduces to cyclotomy.  More precisely, the same resolvent factorization as in \S181 holds, the first auxiliary lemma remains true after the congruence calculation is adjusted for \(2\)-power moduli, and the unit obstruction is removed by replacing the root-of-unity factor by one belonging to a suitable cyclotomic field.  Hence a simple cyclic Abelian field of degree \(2^a\) is a divisor of a cyclotomic field.

Consequently theorem I is established in all cases:

\textbf{Every Abelian field over the rational numbers is contained in a cyclotomic field.}

Equivalently, every finite Abelian extension of \(\mathbb Q\) is generated by roots of unity.  The construction is not merely existential.  The preceding proof shows how the required roots of unity arise from the Lagrange resolvents of a cyclic equation of prime-power degree, and how the remaining ideal and unit factors are removed by the two auxiliary theorems.

\booktitle{Twenty-third Section. Class Number}

\webersection{185}{Dirichlet's Theorem on Units}

The investigations now to be made are quite general; they apply to every algebraic number field and need not be restricted to cyclotomic fields.  They originate in Dirichlet's work on units and class numbers, and were put into the geometric form used here by Minkowski.

Let \(\Omega\) be a field of degree \(n\), and let
\[
  \Omega^{(1)},\Omega^{(2)},\ldots,\Omega^{(n)}
\]
be the conjugate fields.  Some of these are real, while the imaginary ones occur in conjugate pairs.  We count each real embedding once and each pair of conjugate imaginary embeddings once.  Let the resulting number be \(v\).  Thus, if all conjugate fields are real, \(v=n\); if all are imaginary, \(n=2v\).  The number of imaginary pairs is \(n-v\), and the number of real embeddings is \(2v-n\).

Choose one representative from each imaginary pair.  To the resulting \(v\) embeddings attach the numbers
\[
  d_s=\begin{cases}
  1,&\text{if the embedding is real},\\
  2,&\text{if it represents an imaginary pair},
  \end{cases}
\qquad
  \sum_{s=1}^v d_s=n.
\tag{1}
\]
If \(\eta\ne0\) is a number of \(\Omega\), and \(\eta_s\) its conjugate under the \(s\)-th chosen embedding, define its conjugate logarithms by
\[
  \lambda_s=d_s\log|\eta_s|\qquad(s=1,
\ldots,v).
\tag{2}
\]
For an imaginary pair this means
\[
  \lambda_s=\log(\eta_s\overline{\eta_s}).
\]
Therefore
\[
  \lambda_1+\cdots+\lambda_v=\log|
\N(\eta)|.
\tag{3}
\]

If \(\eta\) is a whole number, its absolute norm is a positive integer.  It is equal to \(1\) exactly when \(\eta\) is a unit.  Thus:

\textbf{1.} The sum of the conjugate logarithms of a whole number is positive, and it is zero exactly when the number is a unit.

Let \(\omega_1,
\ldots,
\omega_n\) be a minimal basis of the whole numbers of \(\Omega\).  Every whole number of a conjugate field can be written
\[
  \eta_s=x_1\omega_{1s}+x_2\omega_{2s}+\cdots+x_n\omega_{ns}
\tag{4}
\]
with rational integral \(x_i\).  The logarithmic variables \(\lambda_s\) define regions in the real space of the embeddings.  By choosing bounds for the \(\lambda_s\), and applying the lattice-point principle already used in Minkowski's theorem, one proves the following fundamental assertion.

\textbf{4.} If \(g_1,
\ldots,g_v\) are real numbers not all equal, then there is a unit \(\varepsilon\) of \(\Omega\) such that
\[
  g_1\ell_1+g_2\ell_2+\cdots+g_v\ell_v\ne0,
\tag{5}
\]
where \(\ell_s=d_s\log|\varepsilon_s|\) are the conjugate logarithms of \(\varepsilon\).

The proof proceeds by choosing infinitely many whole numbers whose logarithms lie in a moving strip.  Since their norms remain below a fixed bound, two of them have the same norm and the same residues modulo that norm.  Their quotient is then a unit.  The displacement of their logarithms can be made to have non-zero scalar product with the prescribed vector \((g_s)\).  This is the Dirichlet-Minkowski lemma on which the structure theorem for units rests.

\webersection{186}{Systems of Independent Units and Exponent Systems}

Assume from now on that \(v>1\).  Taking in the preceding theorem, for example, \(g_1=1\) and all other \(g_s=0\), one sees that there is a unit for which at least one conjugate logarithm is not zero.  Since all powers of such a unit have the same property, there are infinitely many units.

The preceding result has a stronger form.  For every integer
\[
  s\le v-1
\]
one can choose units
\[
  \varepsilon_1,\varepsilon_2,
\ldots,\varepsilon_s
\]
with conjugate logarithms
\[
  \ell_{i1},\ell_{i2},\ldots,
\ell_{iv}
\]
so that some \(s\)-rowed determinant formed from the matrix \((\ell_{ik})\) is non-zero.  In particular one can arrange, after a suitable choice of columns, that
\[
  L_s=\det(\ell_{ij})_{1\le i,j\le s}
e0.
\tag{1}
\]
The proof is by induction.  The case \(s=1\) is the result just mentioned.  If units have been chosen for \(s-1\), expand the determinant \(L_s\) by the elements of its last row:
\[
  L_s=g_1\ell_{s1}+g_2\ell_{s2}+\cdots+g_s\ell_{ss}.
\tag{2}
\]
The coefficients \(g_i\) are minors from the already chosen rows, and one of them is non-zero.  Taking the remaining coefficients to be zero in the unused columns, not all \(g_1,
\ldots,g_v\) are equal.  The theorem of \S185 supplies a unit \(\varepsilon_s\) for which (2) is non-zero.

For \(s=v\) such a determinant must vanish.  Indeed for every unit
\[
  \ell_1+\ell_2+\cdots+\ell_v=0,
\tag{3}
\]
so all logarithmic vectors of units lie in the hyperplane (3).  Hence the maximal number of independent unit directions is \(v-1\).

Let
\[
  \varepsilon_1,
\ldots,
\varepsilon_{v-1}
\tag{4}
\]
be a system for which a determinant
\[
  L=L(\varepsilon_1,
\ldots,
\varepsilon_{v-1})
\tag{5}
\]
is non-zero.  If \(\varepsilon\) is another unit, its logarithmic vector can be expressed uniquely as a rational linear combination of the logarithmic vectors of (4):
\[
  \ell(\varepsilon)=\xi_1\ell(\varepsilon_1)+\cdots+\xi_{v-1}\ell(\varepsilon_{v-1}).
\tag{6}
\]
The rational numbers \(\xi_i\) are obtained by Cramer's rule from determinants of logarithms.  They are called the exponents of \(\varepsilon\) with respect to the chosen system of units.  Addition of logarithms shows that multiplication of units corresponds to addition of their exponent systems.

Thus the units form, up to roots of unity, a lattice in the \((v-1)\)-dimensional logarithmic hyperplane.  The next paragraph chooses a fundamental basis of this lattice.

\webersection{187}{Fundamental Systems of Units}

Among all systems of \(v-1\) independent units choose one for which the absolute value of the determinant \(L\) is as small as possible.  This is possible because, after a fixed independent system has been chosen, every other determinant is a rational multiple of the first with a denominator bounded by that system; among non-zero rational multiples with a common denominator there is a least absolute value.  Such a system is called a fundamental system of units, and the minimal positive value of \(|L|\) is the regulator of the field.

Let
\[
  \varepsilon_1,
\varepsilon_2,
\ldots,
\varepsilon_{v-1}
\tag{1}
\]
be a fundamental system.  We prove that every unit has integral exponents with respect to it.  Suppose, on the contrary, that a unit \(\varepsilon\) has an exponent system not consisting entirely of integers.  By multiplying \(\varepsilon\) by suitable integral powers of the \(\varepsilon_i\), one obtains a unit \(\varepsilon_0\) whose exponents are proper fractions, not all zero.  If, say, the first exponent \(\xi_1\) is a non-zero proper fraction, then replacing \(\varepsilon_1\) by \(\varepsilon_0\) changes the determinant by the factor \(\xi_1\):
\[
  L(\varepsilon_0,
\varepsilon_2,
\ldots,
\varepsilon_{v-1})=
  \xi_1 L(\varepsilon_1,
\varepsilon_2,
\ldots,
\varepsilon_{v-1}).
\tag{2}
\]
This contradicts the minimality of \(|L|\).  Hence all exponents are integers.

It remains to know how far a unit is determined by its exponent system.  If two units \(\varepsilon',\varepsilon''\) have the same exponents, then their quotient
\[
  \rho=\varepsilon'/\varepsilon''
\]
has all logarithms equal to zero.  Therefore every conjugate of \(\rho\) has absolute value \(1\).  By Kronecker's theorem, \(\rho\) is a root of unity.  We have proved:

\textbf{I.} In every algebraic number field \(\Omega\) there is a system of \(v-1\) fundamental units
\[
  \varepsilon_1,
\ldots,
\varepsilon_{v-1}
\]
such that every unit of the field is represented once and only once in the form
\[
  \varepsilon=\rho\,
  \varepsilon_1^{a_1}\varepsilon_2^{a_2}\cdots\varepsilon_{v-1}^{a_{v-1}},
\tag{3}
\]
where the exponents \(a_i\) run independently through all rational integers and \(\rho\) runs through the roots of unity contained in \(\Omega\).

Every other fundamental system is obtained from (1) by an integral linear substitution of determinant \(\pm1\) on the exponents.  Conversely, any such substitution gives another fundamental system, since it changes the determinant \(L\) only by the factor \(\pm1\).

If \(v=1\), the theorem says only that the units are roots of unity.  This is the case for the rational field and for imaginary quadratic fields.  In a real quadratic field one has \(v=2\), and the familiar Pell theory is recovered: all units are powers of one fundamental unit, multiplied by \(\pm1\).

\webersection{188}{Reduced Numbers}

The regulator and the logarithmic representation of the units make it possible to choose, in every associate class of whole numbers, representatives subject to fixed inequalities.  These representatives will be called reduced.

Let \(\eta\) be a non-zero whole number and let
\[
  \lambda_s=d_s\log|\eta_s|
\]
be its conjugate logarithms.  Multiplication by a unit \(\varepsilon_1^{a_1}\cdots\varepsilon_{v-1}^{a_{v-1}}\)
adds to the vector \((\lambda_s)\) an integral linear combination of the logarithmic vectors of the fundamental units.  Hence, by choosing the exponents \(a_i\), one may place the resulting logarithmic vector in a fundamental parallelepiped of the unit lattice.

Thus every associate class contains a number whose logarithms satisfy a fixed system of inequalities
\[
  0\le u_i<1\qquad (i=1,
\ldots,v-1),
\tag{1}
\]
when expressed in the unit coordinates.  The last coordinate is then determined by the norm relation
\[
  \lambda_1+\cdots+\lambda_v=
\log|
\N(\eta)|.
\tag{2}
\]
A whole number satisfying these inequalities will be called reduced.

Only finitely many reduced whole numbers have absolute norm below a given bound.  For (1) confines the logarithmic vector to a bounded region once (2) is bounded, and the corresponding embedded lattice has only finitely many points in a bounded region.  Conversely, every non-zero whole number is associated with a reduced one.  Hence counting non-associated whole numbers of bounded norm may be reduced to counting reduced representatives in a fundamental domain.

This preparation is the bridge between the algebra of ideals and the geometry of numbers.  The next paragraph applies it to the set of whole numbers divisible by a fixed ideal.

\webersection{189}{Bounds for the Number of Whole Numbers Divisible by an Ideal}

Let \(\ideala\) be an ideal of \(\Omega\).  We wish to count, for large \(t\), the number \(T\) of non-associated whole numbers divisible by \(\ideala\) whose absolute norm is less than \(t\).  In the rational field this is just the elementary question of counting natural numbers less than \(t\) and divisible by a fixed integer.  In a general field the answer is no longer exact, but the limiting quotient \(T/t\) exists and can be computed.

Choose a basis form of the ideal \(\ideala\):
\[
  \alpha=a_1x_1+a_2x_2+\cdots+a_nx_n,
\tag{1}
\]
so that, when the \(x_i\) run through all rational integers, \(\alpha\) runs through all whole numbers divisible by \(\ideala\).  Let
\[
  y_s=a_{1s}x_1+a_{2s}x_2+\cdots+a_{ns}x_n
\tag{2}
\]
be the conjugate values.  If \(A\) is the determinant of this linear system, then
\[
  A^2=
\N(\ideala)^2\Delta,
\tag{3}
\]
where \(\Delta\) is the discriminant of the field, taken with its sign, and \(\N(\ideala)\) is the absolute norm of the ideal.

The integral systems \((x_1,
\ldots,x_n)\) are lattice points.  To count the numbers \(\alpha\), we consider a region \(S\) defined by the following conditions: \(\alpha\) is reduced in the sense of \S188, and
\[
  |
\N(\alpha)|<t.
\tag{4}
\]
After the scaling
\[
  x_i'=t^{-1/n}x_i,
\]
the region becomes independent of \(t\), up to similarity.  If \(Z_t\) is the number of lattice points in the corresponding enlarged region, then the quotient \(Z_t/t\) tends to the volume of the fixed region.

The reduced numbers are counted with the multiplicity coming from the roots of unity.  If \(w\) denotes the number of roots of unity in \(\Omega\), then each class of associated whole numbers contains exactly \(w\) associates, and therefore
\[
  Z_t=wT.
\tag{5}
\]
Thus the desired limiting quotient is obtained from the volume of the reduced region divided by \(w\).  The explicit computation of that volume is the subject of the next paragraph.

\webersection{190}{Determination of the Volume}

We compute the volume of the region used in the preceding count.  The transformation from the integral coordinates \(x_i\) to the conjugate coordinates \(y_s\) has determinant \(A\).  For each imaginary pair, write
\[
  y=u_1+iu_2,
\qquad \overline y=u_1-iu_2.
\tag{1}
\]
The real coordinates consist of all real conjugates and the real and imaginary parts of one member of each imaginary pair.  Passing from the complex pairs to these real coordinates contributes the factor \(2^{-(n-v)}\) in absolute value, and the discriminant is taken in absolute value.  Hence the Jacobian ultimately contains
\[
  \frac{2^{n-v}}{\N(\ideala)\sqrt{|\Delta|}}.
\tag{2}
\]

Next introduce polar coordinates for each imaginary pair:
\[
  u_1=e^{\lambda/2}\cos\theta,
\qquad
  u_2=e^{\lambda/2}\sin\theta,
\qquad 0\le\theta<2\pi.
\tag{3}
\]
For a real conjugate, write
\[
  u=\pm e^\lambda.
\tag{4}
\]
The signs of the real conjugates split the region into \(2^{2v-n}\) equal parts.  The angles from the imaginary pairs give factors \(\pi\), and the logarithmic variables are then changed to the unit coordinates
\[
  \lambda_s=\xi_1\ell_{1s}+\cdots+\xi_{v-1}\ell_{v-1,s}+\frac{1}{v}\log|
\N(\alpha)|.
\tag{5}
\]
The determinant of this change of variables is, up to sign, the regulator \(L\), multiplied by the number \(w\) of roots of unity when associated elements are identified.

Carrying out the integration gives Weber's constant
\[
  g=\frac{2^{2v-n}\pi^{n-v}L}{w\sqrt{|\Delta|}}.
\tag{6}
\]
Here \(L\) is the regulator of the field, \(w\) the number of roots of unity, \(n-v\) the number of imaginary pairs, and \(2v-n\) the number of real embeddings.

The result is:

\textbf{1.} Let \(T\) be the number of non-associated whole numbers divisible by \(\ideala\) whose absolute norm is smaller than \(t\).  Then
\[
  \lim_{t\to\infty}\frac{T}{t}
  =\frac{g}{\N(\ideala)}.
\tag{7}
\]
The constant \(g\) is positive and depends only on the field \(\Omega\), not on the ideal \(\ideala\).

Formula (7) is the asymptotic form of the geometry-of-numbers calculation.  It is the analytic input from which the class number formula follows.

\webersection{191}{Theorems from the Theory of Series}

A few elementary facts about infinite series are needed for the application to ideal classes.

\textbf{1.} Let \(a_1,a_2,
\ldots\) be real numbers, positive, negative, or zero, and suppose that all partial sums
\[
  A_n=a_1+
\cdots+a_n
\]
remain bounded in absolute value by a fixed constant.  Then, for every positive \(s\), the series
\[
  \sum_{n=1}^\infty \frac{a_n}{n^s}
\tag{1}
\]
converges, and its sum is a continuous function of \(s\) for \(s>0\).  This is summation by parts: the tail is controlled by the bounded partial sums and by the monotonic decrease of \(n^{-s}\).

\textbf{2.} Let \(\mu_1,\mu_2,
\ldots\) be positive numbers ordered increasingly, and suppose there are positive constants \(\alpha,\beta\) such that, for all sufficiently large \(n\),
\[
  \alpha<\frac{n}{\mu_n}<\beta.
\tag{2}
\]
Then
\[
  \sum_{n=1}^{\infty}\frac1{\mu_n^s}
\tag{3}
\]
converges for \(s>1\), and the product
\[
  (s-1)
  \sum_{n=1}^{\infty}\frac1{\mu_n^s}
\tag{4}
\]
remains between bounds arbitrarily close to \(\alpha\) and \(\beta\) as \(s\) approaches \(1\) from above.

The proof compares the sum with the integral
\[
  \int_m^\infty x^{-s}\,dx=\frac{m^{1-s}}{s-1}.
\]
If the limiting quotient
\[
  \lim_{n\to\infty}\frac{n}{\mu_n}=\gamma
\tag{5}
\]
exists, then the sharper conclusion is
\[
  \lim_{s\downarrow1}(s-1)
  \sum_{n=1}^{\infty}\frac1{\mu_n^s}=\gamma.
\tag{6}
\]

Equivalently, let \(T(t)\) be the number of \(\mu_n\) not exceeding \(t\).  If either of the limits
\[
  \lim_{n\to\infty}\frac{n}{\mu_n},
  \qquad
  \lim_{t\to\infty}\frac{T(t)}{t}
\tag{7}
\]
exists and is finite, then the other exists and has the same value.  This is the elementary relation between an ordered sequence and its counting function.

Finally, Weber records the familiar expansion of the Riemann series near \(s=1\).  The series
\[
  \zeta(s)=\sum_{n=1}^{\infty}\frac1{n^s}
\]
has a simple pole at \(s=1\), and
\[
  \lim_{s\downarrow1}(s-1)\zeta(s)=1.
\tag{8}
\]
Moreover the difference
\[
  \zeta(s)-\frac1{s-1}
\]
has a finite limit as \(s\to1\).  This follows from the standard integral representation with the Gamma function,
\[
  \frac1{n^s}=\frac1{\Gamma(s)}\int_0^\infty x^{s-1}e^{-nx}\,dx,
\]
and the resulting integral
\[
  \frac1{\Gamma(s)}\int_0^\infty x^{s-1}
  \left(\frac1{e^x-1}-\frac1x\right)
  \,dx.
\]
At \(s=1\) this gives Euler's constant.  Only the finiteness of the limit is needed below.

\webersection{192}{Application to the Determination of the Class Number}

We now apply these theorems to the ideals of the field \(\Omega\).  Let \(\mu_n\) be the norms of all ideals of \(\Omega\), ordered with repetition and in non-decreasing order.  Let \(T(t)\) be the number of ideals whose norm does not exceed \(t\).  The series
\[
  \Phi(s)=\sum_{\ideala}\frac1{\N(\ideala)^s}
\tag{1}
\]
then converges for \(s>1\).

The ideals split into finitely many ideal classes,
\[
  A_1,A_2,
\ldots,A_h,
\]
where \(h\) is the class number.  Choose for each class \(A_i\) a whole functional \(\varphi_i\) belonging to the inverse class.  Then \(\varphi_i\ideala\) is principal for every \(\ideala\in A_i\), say
\[
  \varphi_i\ideala=(\alpha).
\tag{2}
\]
If \(\N(\ideala)
\le t\), then
\[
  |
\N(\alpha)|\le \N(\varphi_i)t.
\tag{3}
\]
Thus the number of ideals in the class \(A_i\) of norm at most \(t\) is the number of non-associated whole numbers divisible by \(\varphi_i\) and of norm at most \(\N(\varphi_i)t\).  By the asymptotic formula of \S190,
\[
  \lim_{t\to\infty}\frac{T_i(t)}{t}=g.
\tag{4}
\]
The constant is the same for every class.  Summing over the \(h\) classes gives
\[
  \lim_{t\to\infty}\frac{T(t)}{t}=gh.
\tag{5}
\]
By the theorem on series, this is equivalent to
\[
  gh=\lim_{s\downarrow1}(s-1)
  \sum_{\ideala}\frac1{\N(\ideala)^s}.
\tag{6}
\]

The series (1) can be transformed into an Euler product.  If \(\idealp\) runs over all prime ideals, then
\[
  1+\N(\idealp)^{-s}+\N(\idealp)^{-2s}+\cdots
  =\frac1{1-\N(\idealp)^{-s}}.
\]
Unique factorization of ideals gives
\[
  \Phi(s)=\prod_{\idealp}
  \frac1{1-\N(\idealp)^{-s}}.
\tag{7}
\]
If a rational prime \(p\) splits in \(\Omega\) into prime ideals of degrees
\[
  f_1,
\ldots,f_e,
\]
then its contribution is
\[
  \prod_{j=1}^{e}\frac1{1-p^{-f_j s}}.
\tag{8}
\]
Formula (6), together with (7), is the general analytic class number formula in Weber's form:
\[
  gh=\lim_{s\downarrow1}(s-1)
  \prod_{\idealp}
  \frac1{1-\N(\idealp)^{-s}}.
\tag{9}
\]

A useful consequence concerns prime ideals of first degree.  In the product (7), the partial product over prime ideals of degree greater than one remains finite and non-zero even for \(s=1\), because it is dominated by products built from \(p^{-2s}\).  Therefore the pole at \(s=1\) must come from the first-degree prime ideals.  It follows that the product
\[
  \prod_{\N(\idealp)=p}
  (1-p^{-s})
\tag{10}
\]
formed over prime ideals of first degree has, as \(s\to1\), a finite non-zero limiting behavior after the necessary factor has been separated.

In particular:

\textbf{1.} In every algebraic number field there are infinitely many prime ideals of first degree.

For if there were only finitely many, the first-degree part of the Euler product would be non-zero at \(s=1\), and the remaining factors would also be finite and non-zero, contradicting the fact that \(\Phi(s)\) has a simple pole there with residue \(gh>0\).

This theorem will be applied in the next section to cyclotomic fields, where the degrees of prime ideals are described by congruence classes modulo the conductor.

\booktitle{Twenty-fourth Section. Class Number of Cyclotomic Fields}

\webersection{193}{Class Number Representation in the Cyclotomic Field \(\Om\)}

We now apply the general class number formula to the full cyclotomic field \(\Om\), still under the assumption
\[
  m=q^a,\qquad q\ne2,\qquad\mu=\varphi(m)=(q-1)q^{a-1}.
\tag{1}
\]
In \(\Om\), the prime \(q\) has a single prime divisor of degree \(1\), repeated with ramification index \(\mu\).  A rational prime \(p\ne q\), belonging to the exponent \(f\) modulo \(m\), decomposes into
\[
  e=\mu/f
\]
prime ideals of degree \(f\).  Thus the Euler product of \S192 becomes
\[
  \Phi(s)=\frac1{1-q^{-s}}
  \prod_{p\ne q}
  \frac1{(1-p^{-fs})^e}.
\tag{2}
\]
Here \(f=f(p)\) is the least positive exponent for which
\[
  p^f\equiv1\pmod m.
\tag{3}
\]

To transform (2), use the characters of the Abelian group of reduced residue classes modulo \(m\).  Denote this group by \(G\); its order is \(\mu\).  If \(\chi\) runs through all characters of \(G\), then for a prime \(p\ne q\) of exponent \(f\), the values \(\chi(p)\) run through all \(f\)-th roots of unity, each repeated \(e\) times.  Hence
\[
  \prod_\chi (1-\chi(p)x)=(1-x^f)^e.
\tag{4}
\]
Putting \(x=p^{-s}\) gives
\[
  \Phi(s)=
  \prod_{\chi}\left(\sum_{(n,q)=1}\frac{\chi(n)}{n^s}\right).
\tag{5}
\]
The product extends over all \(\mu\) characters modulo \(m\).

The factor belonging to the principal character is
\[
  \sum_{(n,q)=1}\frac1{n^s}
  =\left(1-q^{-s}\right)\zeta(s),
\tag{6}
\]
and therefore
\[
  \lim_{s\downarrow1}(s-1)
  \sum_{(n,q)=1}\frac1{n^s}=1-\frac1q.
\tag{7}
\]
All other character series are continuous at \(s=1\).  Indeed the partial sums of a non-principal character over complete residue systems modulo \(m\) vanish, and so the theorem of \S191 applies.

Consequently, if
\[
  X(\chi)=\sum_{(n,q)=1}\frac{\chi(n)}{n}
\tag{8}
\]
for non-principal characters, the class number formula takes the form
\[
  gh=\left(1-\frac1q\right)
  \prod_{\chi\ne1} X(\chi).
\tag{9}
\]
This is Weber's first class number representation for the cyclotomic field.  The constant \(g\) is the field constant of \S190, expressed in terms of the regulator, discriminant, roots of unity, and number of real and imaginary embeddings of \(\Om\).

\webersection{194}{Determination of the Sums \(X\)}

It remains to express the sums
\[
  X(\chi)=\sum_{(n,q)=1}\frac{\chi(n)}{n}
\tag{1}
\]
by finite formulae.  Since \(\chi(n)\) is periodic modulo \(m\), write \(t\) for the least positive residue of \(n\) and decompose
\[
  n=t+ml.
\]
Then
\[
  \chi(n)=\chi(t).
\]
Using
\[
  \frac1n=\int_0^1 x^{n-1}\,dx,
\]
we obtain
\[
  X(\chi)=\int_0^1
  \sum_{t=1}^{m}\chi(t)x^{t-1}
  \left(1+x^m+x^{2m}+\cdots\right)
  dx.
\tag{2}
\]
Thus
\[
  X(\chi)=\int_0^1
  \frac{f(x)}{x(1-x^m)}\,dx,
\tag{3}
\]
where
\[
  f(x)=\sum_{t=1}^{m}\chi(t)x^t.
\tag{4}
\]
For a non-principal character,
\[
  f(0)=0,
  \qquad
  f(1)=\sum_{t=1}^{m}\chi(t)=0,
\tag{5}
\]
so \(f(x)\) is divisible by \(x(1-x)\).  Decompose the rational function in (3) into partial fractions with respect to the roots \(\zetaS^s\) of \(x^m-1\).  One obtains
\[
  \frac{f(x)}{x(1-x^m)}
  =\frac1m\sum_{s=1}^{m-1}
  \frac{f(\zetaS^s)}{x-\zetaS^s}.
\tag{6}
\]
The elementary integral
\[
  \int_0^1\frac{dx}{x-e^{i\theta}}
\]
may be expressed through logarithms and angles; in the resulting formula the terms combine according to the value of \(\chi(-1)\).  Since
\[
  f(\zetaS^{-s})=\chi(-1)f(\zetaS^s),
\tag{7}
\]
the characters divide into two sorts:
\[
  \chi_1(-1)=+1,
  \qquad
  \chi_2(-1)=-1.
\tag{8}
\]
For even characters \(\chi_1\), the paired terms give logarithmic expressions; for odd characters \(\chi_2\), they give trigonometric cotangent expressions.  Weber writes the resulting finite formulae in terms of
\[
  \log\left(4\sin^2\frac{\pi s}{m}\right)
\]
for the even characters and in terms of the finite sums involving \(\zetaS^s\) for the odd characters.  In compact form the two cases are:
\[
  X(\chi_1)=-\frac1m
  \sum_{s=1}^{m-1}f_{\chi_1}(\zetaS^s)
  \log\left(4\sin^2\frac{\pi s}{m}\right),
\tag{9}
\]
and
\[
  X(\chi_2)=-\frac{\pi i}{m}
  \sum_{s=1}^{m-1} f_{\chi_2}(\zetaS^s)
  \frac{1+\zetaS^s}{1-\zetaS^s}.
\tag{10}
\]
These formulae replace the infinite series (1) by finite expressions in roots of unity and logarithms of cyclotomic sine factors.  They are the computational form needed for the class number formula of \S193 and prepare the separation, in the next paragraph, of the class number of the real subfield \(\Hm\) from the remaining factor of the full cyclotomic class number.

\end{document}
